\documentclass{article}

\usepackage{hyperref}

\newcommand{\meetingdate}{8 July 2026}
\usepackage[
    a4paper,
    top=3.5cm,
    bottom=2.3cm,
    left=2.3cm,
    right=2.3cm,
    headheight=42pt
]{geometry}
\usepackage[utf8]{inputenc}
\usepackage[T1]{fontenc}
\usepackage{fancyhdr}
\usepackage{lastpage}
\usepackage{enumitem}
\usepackage{amssymb}
\usepackage{etoolbox}
\usepackage{xcolor}
\usepackage{tikz}
\usepackage{tcolorbox}

\newcommand{\project}{Interactive VR Museum}
\newcommand{\meetingtype}{Weekly Meeting}
\providecommand{\meetingdate}{\today}

\newcommand{\authorname}{Jannick Seper}
\newcommand{\tutors}{Rahel Arnold, Raphael Waltenspül}

\lhead{\project\\\meetingtype\\Date: \meetingdate}
\rhead{\authorname\\\tutors\\Page \thepage\ of \pageref{LastPage}}
\pagestyle{fancy}

\newcommand{\done}{
    \section*{1. What I have done}
}

\newcommand{\support}{
    \section*{2. Where I need support}
}

\newcommand{\plan}{
    \section*{3. What I am planning to do for the next meeting}
}

\newtcolorbox{checkpointbox}{
    colback=gray!5,
    colframe=gray!50,
    boxrule=0.5pt,
    arc=3mm,
    left=3mm,
    right=3mm,
    top=2mm,
    bottom=2mm,
    width=\textwidth
}

\newcommand{\checkpoint}[2][open]{
    \begin{checkpointbox}
        \ifstrequal{#1}{done}
            {$\boxtimes$}
            {$\square$}
        \hspace{0.3cm} #2
    \end{checkpointbox}
}

\begin{document}

\done
\begin{itemize}
    \item Implemented similarity search and made it accessible from the museum through a popup-based workflow.
    \item Added vector values from vitrivr to displayed media nodes so that similarity searches can be started from existing results.
    \item Refactored the Godot-side code and backend integration.
    \item Added hand tracking support for Meta Quest 3, including switching between controller and hand tracking modes for on-device usage and streaming. Some errors still remain.
    \item Improved visual quality and reduced warnings, with the tradeoff of lower performance.
    \item Spent significant time investigating how to export a runnable Godot build for Apple Vision Pro.
    \item Created first evaluation points and uploaded the German evaluation concept.
    \item Added templates for the midterm presentation and the written report.
\end{itemize}

\support
\begin{itemize}
    \item Access to a Vive XR Elite headset for testing.
    \item Access to more videos and 3D assets.
    \item Feedback on the planned evaluation points.
    \item .NET xrOS integration is not available yet, so we need to find a way to test the Apple Vision Pro. Maybe streaming over ALVR?
\end{itemize}

\plan
\checkpoint{Test streaming to the Apple Vision Pro}
\checkpoint{Create a new room for original-size 3D assets}
\checkpoint{Test hand tracking on additional devices}
\checkpoint{Fix remaining hand tracking and input mode issues}
\checkpoint{Prepare for the midterm presentation}

\end{document}
